% SGA 3, Exposé IX, tradução brasileira.
% Escopo: §2 integral.
% Autoridade: Polo--Gille Exp9-8nov09.pdf, pp. locais 4--7,
% leitor completo pp. 650--653, pp. impressas 41--46.

\setcounter{footnote}{2}

\SGARefTarget{sga3:ix:section:2}\section*{2. Extensão de certas
propriedades dos grupos diagonalizáveis aos grupos de tipo
multiplicativo}
\addcontentsline{toc}{section}{2. Extensão de certas propriedades dos
grupos diagonalizáveis aos grupos de tipo multiplicativo}

As extensões em questão são consequências essencialmente triviais
dos resultados da Exposição precedente, levando-se em conta
\SGARefLink{sga3:ix:definition:1:1}{as definições de 1.1} e o caráter
«local» (para a topologia fielmente plana e quase compacta) dos
resultados considerados.

\medskip
\noindent\SGARefTarget{sga3:ix:definition:2:0}\textbf{Definição 2.0.}
\footnote{Nota do editor: acrescentou-se esta definição; ela intervém nas
Proposições~\SGARefLink{sga3:ix:proposition:2:3}{2.3} e
\SGARefLink{sga3:ix:proposition:2:7}{2.7}.}
\textit{Denotamos por \(\mathcal M\) o conjunto dos morfismos
fielmente planos e quase compactos.}

\medskip
\noindent\SGARefTarget{sga3:ix:proposition:2:1}\textbf{Proposição 2.1.}
\textit{Seja \(G\) um grupo de tipo multiplicativo sobre um pré-esquema
\(S\). Então:}
\begin{enumerate}
  \item[(a)]\SGARefTarget{sga3:ix:proposition:2:1:item:a:01}
  \textit{\(G\) é fielmente plano sobre \(S\) e afim sobre \(S\)
  (\emph{a fortiori}, quase compacto sobre \(S\));}

  \item[(b)]\SGARefTarget{sga3:ix:proposition:2:1:item:b:01}
  \textit{\(G\) é de tipo finito sobre \(S\) se, e somente se, é de
  apresentação finita sobre \(S\), se, e somente se, para todo \(s\in S\),
  o tipo de \(G\) em \(s\) é dado por um grupo comutativo de tipo
  finito;}

  \item[(c)]\SGARefTarget{sga3:ix:proposition:2:1:item:c:01}
  \textit{\(G\) é finito sobre \(S\) se, e somente se, para todo
  \(s\in S\), o tipo de \(G\) em \(s\) é dado por um grupo
  comutativo finito, se, e somente se (quando \(S\) é quase compacto), \(G\)
  é de tipo finito sobre \(S\) e é anulado por um inteiro \(n>0\);}

  \item[\((c')\)]\SGARefTarget{sga3:ix:proposition:2:1:item:cprime:01}
  \textit{\(G\) é inteiro sobre \(S\) se, e somente se, para todo
  \(s\in S\), o tipo de \(G\) em \(s\) é dado por um grupo
  comutativo de torção;}

  \item[(d)]\SGARefTarget{sga3:ix:proposition:2:1:item:d:01}
  \textit{\(G\) é o grupo unidade se, e somente se, para todo \(s\in S\),
  o tipo de \(G\) em \(s\) é dado pelo grupo comutativo zero;}

  \item[(e)]\SGARefTarget{sga3:ix:proposition:2:1:item:e:01}
  \textit{\(G\) é liso sobre \(S\) se, e somente se, para todo \(s\in S\),
  o tipo de \(G\) em \(s\) é dado por um grupo comutativo de tipo
  finito cujo subgrupo de torção tem ordem prima com a característica
  de \(\kappa(s)\).}
\end{enumerate}

Essas afirmações deduzem-se de
\SGARefLink{sga3:viii:proposition:2:1}{VIII~2.1}, levando em conta que
as propriedades em questão descendem por morfismos fielmente planos e
quase compactos (cf. SGA~1, VIII, ou EGA~IV\textsubscript{2}, \S2).

Utilizando \SGARefLink{sga3:viii:corollary:3:5}{VIII~3.5}, obtém-se
também:

\medskip
\noindent\SGARefTarget{sga3:ix:proposition:2:2}\textbf{Proposição 2.2.}
\textit{Seja \(G\) um grupo de tipo multiplicativo e de tipo finito sobre
\(S\). Para todo inteiro \(n\ne0\), o núcleo
\({}_nG=\ker(n\cdot\operatorname{id}_G)\) é um grupo de tipo
multiplicativo, finito sobre \(S\).}

\medskip
\noindent\SGARefTarget{sga3:ix:proposition:2:3}\textbf{Proposição 2.3.}
\textit{Seja \(G\) um grupo de tipo multiplicativo sobre o pré-esquema
\(S\), que atua livremente à direita sobre um \(S\)-pré-esquema
\(X\), afim sobre \(S\). Então:}
\begin{enumerate}
  \item[(a)]\SGARefTarget{sga3:ix:proposition:2:3:item:a:01}
  \textit{a relação de equivalência definida por \(G\) em \(X\) é
  \(\mathcal M\)-efetiva
  (\SGARefLink{sga3:iv:subsection:3:4}{IV~3.4}), e \(Y=X/G\) é afim
  sobre \(S\);}

  \item[(b)]\SGARefTarget{sga3:ix:proposition:2:3:item:b:01}
  \textit{se, além disso, \(X\) é de apresentação finita (respectivamente,
  de tipo finito) sobre \(S\), o mesmo vale para \(Y\).}
\end{enumerate}

\SGARefLink{sga3:ix:proposition:2:3:item:a:01}{A primeira afirmação}
deduz-se de \SGARefLink{sga3:viii:theorem:5:1}{VIII~5.1}, que trata do caso
em que \(G\) é diagonalizável, e de
\SGARefLink{sga3:iv:proposition:3:5:2}{IV~3.5.2}, que permite reduzir-se a
esse caso. Com efeito, os morfismos fielmente planos e quase compactos
\(S'\to S\) são morfismos de descida efetiva para a categoria fibrada
dos esquemas afins sobre os esquemas: se \(Y'\) é afim sobre \(S'\)
e está munido de um dado de descida relativo a \(S'\to S\), esse dado
é efetivo, isto é, \(Y'\) provém de um esquema \(Y\) afim sobre
\(S\) (cf. SGA~1, \SGARefLink{sga3:viii:proposition:2:1}{VIII~2.1}).

% Página-fonte: Exp9 p. local 5; leitor completo p. 651; p. impressa 43.
Para \SGARefLink{sga3:ix:proposition:2:3:item:b:01}{a segunda afirmação},
reduz-se igualmente ao caso diagonalizável de
\SGARefLink{sga3:viii:corollary:5:8}{VIII~5.8}, pois as condições de
finitude consideradas descendem por morfismos fielmente planos e
quase compactos (SGA~1, VIII~3.3 e 3.6\footnote{Nota do editor: veja-se
também V, 9.1, ou EGA~IV\textsubscript{2}, 2.7.1.}). Procedendo como em
VIII, Corolários~5.5--5.7, deduz-se de
\SGARefLink{sga3:ix:proposition:2:3}{2.3}:

\medskip
\noindent\SGARefTarget{sga3:ix:corollary:2:4}\textbf{Corolário 2.4.}
\textit{Sob as hipóteses de
\SGARefLink{sga3:ix:proposition:2:3}{2.3}, o morfismo gráfico}
\SGARefTarget{sga3:ix:corollary:2:4:diagram:01}
\[
  X\times_S G\longrightarrow X\times_S X
\]
\textit{é uma imersão fechada. Para toda seção \(\sigma\) de \(X\)
sobre \(S\), o morfismo correspondente}
\[
  G\longrightarrow X,\qquad g\longmapsto\sigma\cdot g,
\]
\textit{é uma imersão fechada.}

Em particular:

\medskip
\noindent\SGARefTarget{sga3:ix:corollary:2:5}\textbf{Corolário 2.5.}
\textit{Seja \(u:G\to H\) um monomorfismo de \(S\)-pré-esquemas em grupos,
com \(G\) de tipo multiplicativo e \(H\) afim sobre \(S\). Então
\(u\) é uma imersão fechada, \(H/G=Y\) existe e é afim sobre \(S\),
e \(H\) é um fibrado principal homogêneo sobre \(Y\) sob \(G_Y\).}

\medskip
\noindent\SGARefTarget{sga3:ix:remark:2:6}\textbf{Observação 2.6.}
Seja \(u:G\to H\) um monomorfismo de \(S\)-pré-esquemas em grupos, onde
\(G\) é de tipo multiplicativo e de tipo finito sobre \(S\), e \(H\) é
de apresentação finita e separado sobre \(S\). Pode-se então demonstrar
que \(u\) é uma imersão fechada utilizando
\SGARefLink{sga3:viii:proposition:7:12}{VIII~7.12} (veja-se a Observação
\SGARefLink{sga3:viii:remark:7-13:item:b}{VIII~7.13(b)}).

\medskip
\noindent\SGARefTarget{sga3:ix:proposition:2:7}\textbf{Proposição 2.7.}
\textit{Seja \(u:G\to H\) um homomorfismo de \(S\)-grupos de tipo
multiplicativo, com \(H\) de tipo finito sobre \(S\). Então:}
\begin{enumerate}
  \item[(i)]\SGARefTarget{sga3:ix:proposition:2:7:item:i:01}
  \textit{\(G_0=\ker u\) é um \(S\)-grupo de tipo multiplicativo; a
  relação de equivalência definida por \(G_0\) em \(G\) é
  \(\mathcal M\)-efetiva e, portanto
  (\SGARefLink{sga3:iv:subsection:3:4}{IV~3.4}), \(u\) se fatora como}
  \SGARefTarget{sga3:ix:proposition:2:7:diagram:01}
  \[
    G\longrightarrow G/G_0=I\longrightarrow H,
  \]
  \textit{onde \(I\to H\) é um monomorfismo\footnote{Nota do editor:
  na realidade, até mesmo uma imersão fechada.} de \(S\)-grupos. O grupo
  \(I\) é de tipo multiplicativo, a relação de equivalência em \(H\)
  definida por \(I\) é \(\mathcal M\)-efetiva e, por conseguinte,
  \(H_0=H/I\) existe e é também de tipo multiplicativo;}

  \item[(ii)]\SGARefTarget{sga3:ix:proposition:2:7:item:ii:01}
  \textit{\(H_0\) e \(I\) são de tipo finito, e \(G_0\) também o é
  quando \(G\) o é.}
\end{enumerate}

\emph{Demonstração.}
Procedendo como para \SGARefLink{sga3:ix:proposition:2:3}{2.3}, reduz-se
ao caso em que \(G\) e \(H\) são diagonalizáveis; então
\SGARefLink{sga3:ix:proposition:2:7}{2.7} reduz-se a
\SGARefLink{sga3:viii:corollary:3:4}{VIII~3.4}.

Registremos os corolários seguintes.

% Página-fonte: Exp9 p. local 6; leitor completo p. 652; p. impressa 44.
\medskip
\noindent\SGARefTarget{sga3:ix:corollary:2:8}\textbf{Corolário 2.8.}
\begin{enumerate}
  \item[(a)]\SGARefTarget{sga3:ix:corollary:2:8:item:a:01}
  \textit{Seja \(S\) um pré-esquema. A categoria dos \(S\)-grupos de
  tipo multiplicativo e de tipo finito é uma categoria abeliana.}

  \item[(b)]\SGARefTarget{sga3:ix:corollary:2:8:item:b:01}
  \textit{Seja \(u:G\to H\) um homomorfismo nessa categoria. Ele é um
  monomorfismo (respectivamente, um epimorfismo, respectivamente, um
  isomorfismo) nessa categoria se, e somente se, \(u\) é um monomorfismo de
  pré-esquemas (respectivamente, \(u\) é fielmente plano,
  respectivamente, \(u\) é um isomorfismo de pré-esquemas).}
\end{enumerate}

Basta observar que o produto de dois \(S\)-grupos de tipo multiplicativo
ainda é de tipo multiplicativo (o que é imediato), e evidentemente
é de tipo finito sobre \(S\) se ambos os fatores o são. O restante de
\SGARefLink{sga3:ix:corollary:2:8}{2.8} deduz-se imediatamente de
\SGARefLink{sga3:ix:proposition:2:7}{2.7}; deixam-se os pormenores ao leitor.

\medskip
\noindent\SGARefTarget{sga3:ix:corollary:2:9}\textbf{Corolário 2.9.}
\textit{Seja \(u:G\to H\) um homomorfismo de \(S\)-grupos de tipo
multiplicativo e de tipo finito. Seja \(U\) o conjunto dos pontos \(s\in S\)
para os quais \(u_s:G_s\to H_s\) é um monomorfismo (respectivamente,
fielmente plano, respectivamente, um isomorfismo). Então \(U\) é ao mesmo
tempo aberto e fechado, e o homomorfismo induzido \(u_U:G_U\to H_U\) é um
monomorfismo (respectivamente, fielmente plano, respectivamente, um
isomorfismo).}

Seja \(P\) (respectivamente, \(Q\)) o núcleo (respectivamente, o conúcleo)
de \(u\). Por \SGARefLink{sga3:ix:proposition:2:7}{2.7}, \(Q\) existe e
\(P\) e \(Q\) são de tipo multiplicativo; além disso, a formação de \(P\) e
\(Q\) comuta com toda mudança de base \(S'\to S\), em particular com a
formação das fibras. Por outro lado, \(u\) é um monomorfismo
(respectivamente, fielmente plano, respectivamente, um isomorfismo) se, e
somente se, \(P=0\) (respectivamente, \(Q=0\), respectivamente, \(P=Q=0\)).
Essas observações reduzem-nos a verificar a afirmação seguinte: se
\(R\) é um \(S\)-grupo de tipo multiplicativo, o conjunto \(U\) dos
pontos \(s\in S\) para os quais \(R_s=0\) é aberto e fechado, e
\(R_U=0\). Isso está contido na Observação
\SGARefLink{sga3:ix:remark:1:4:1:item:a}{1.4.1(a)}.

\medskip
\noindent\SGARefTarget{sga3:ix:corollary:2:10}\textbf{Corolário 2.10.}
\textit{Seja \(G\) um \(S\)-grupo de tipo multiplicativo e de tipo finito,
e sejam \(H\) e \(H'\) dois subgrupos de tipo multiplicativo.}
\begin{enumerate}
  \item[(a)]\SGARefTarget{sga3:ix:corollary:2:10:item:a:01}
  \textit{\(H''=H\cap H'\) é um subgrupo de \(G\) de tipo
  multiplicativo.}

  \item[(b)]\SGARefTarget{sga3:ix:corollary:2:10:item:b:01}
  \textit{Seja \(U\) o conjunto dos \(s\in S\) para os quais
  \(H_s\subseteq H'_s\) (respectivamente, \(H_s=H'_s\)). Então \(U\)
  é aberto e fechado, e \(H_U\subseteq H'_U\) (respectivamente,
  \(H_U=H'_U\)).}
\end{enumerate}

Aqui, naturalmente, o sinal de interseção em \(H\cap H'\) denota a
interseção em sentido funtorial, a saber, \(H\times_GH'\), que é
manifestamente um \(S\)-subgrupo de \(G\). Do mesmo modo, o sinal de inclusão
(respectivamente, o sinal de igualdade) denota a relação de ordem
(respectivamente, a igualdade) entre subfuntores de \(H\), e não a
inclusão (respectivamente, a igualdade) dos conjuntos subjacentes.

Aplicando primeiro \SGARefLink{sga3:ix:proposition:2:7}{2.7} ao morfismo de
inclusão \(H\to G\), obtém-se que \(H\) é de tipo finito; o mesmo
vale para \(H'\). De
\SGARefLink{sga3:ix:corollary:2:8}{2.8} segue então que
\(H\cap H'\) é de tipo multiplicativo e de tipo finito. (Observe-se que o
funtor canônico da categoria considerada em
\SGARefLink{sga3:ix:corollary:2:8}{2.8} para a categoria
\((\mathbf{Sch})/S\) comuta com os limites projetivos finitos.)

A formação de \(H\cap H'\) comuta com a extensão da base, em
particular com a formação das fibras. Além disso, \(H\subseteq H'\)
(respectivamente, \(H=H'\)) equivale a \(H''=H\) (respectivamente, a
\(H''=H\) e \(H''=H'\)). Consideremos os homomorfismos canônicos
\(H''\to H\) e \(H''\to H'\). Então \(U\) é o conjunto dos pontos
\(s\in S\) para os quais o homomorfismo induzido \(H''_s\to H_s\) é um
isomorfismo (respectivamente, para os quais \(H''_s\to H_s\) e
\(H''_s\to H'_s\) são isomorfismos).

% Página-fonte: Exp9 p. local 7; leitor completo p. 653; pp. impressas 45--46.
Por \SGARefLink{sga3:ix:corollary:2:9}{2.9}, segue que \(U\) é aberto
e fechado e que \(H''_U\to H_U\) (respectivamente, \(H''_U\to H_U\) e
\(H''_U\to H'_U\)) são isomorfismos. Isso dá a conclusão desejada.

\medskip
\noindent\SGARefTarget{sga3:ix:proposition:2:11}\textbf{Proposição 2.11.}
\textit{Seja \(S\) um pré-esquema, seja \(G\) um \(S\)-grupo de tipo
multiplicativo e de tipo finito sobre \(S\), seja \(H\) um subgrupo de
\(G\) de tipo multiplicativo, e ponhamos \(K=G/H\), grupo quociente de
\(G\) de tipo multiplicativo.}
\begin{enumerate}
  \item[(i)]\SGARefTarget{sga3:ix:proposition:2:11:item:i:01}
  \textit{Suponhamos que \(G\) seja trivial, isto é, diagonalizável.
  Existe uma partição de \(S\) em subconjuntos abertos e fechados \(S_i\)
  tal que, para todo \(i\), \(H_{S_i}\) e \(K_{S_i}\) são
  diagonalizáveis. Em particular, se \(S\) é conexo, \(H\) e \(K\) são
  diagonalizáveis.}

  \item[(ii)]\SGARefTarget{sga3:ix:proposition:2:11:item:ii:01}
  \textit{\SGARefLink{sga3:ix:proposition:2:11:item:i:01}{A mesma
  afirmação} vale substituindo «diagonalizável» por «isotrivial», desde
  que \(S\) seja conexo, localmente conexo ou quase compacto.}

  \item[(iii)]\SGARefTarget{sga3:ix:proposition:2:11:item:iii:01}
  \textit{Se \(G\) é localmente trivial (respectivamente, localmente
  isotrivial, respectivamente, quase-isotrivial), o mesmo vale para \(K\)
  e \(H\).}
\end{enumerate}

\emph{Demonstração.}
\SGARefLink{sga3:ix:proposition:2:11:item:i:01}{(i)} Por hipótese,
\(G=D_S(M)\), onde \(M\) é um grupo comutativo de tipo finito. Para
cada grupo quociente \(M_i\) de \(M\), seja \(H_i=D_S(M_i)\) o subgrupo
diagonalizável correspondente de \(G\)
(\SGARefLink{sga3:viii:theorem:3:1}{VIII~3.1}). Seja \(S_i\) o conjunto
dos pontos \(s\in S\) para os quais \(H_s=(H_i)_s\). Por
\SGARefLink{sga3:ix:corollary:2:10}{2.10}, \(S_i\) é aberto e fechado,
e \(H_{S_i}=(H_i)_{S_i}\); assim, \(H_{S_i}\) é diagonalizável e, por
conseguinte, \(K_{S_i}\) também o é
(\SGARefLink{sga3:viii:theorem:3:1}{VIII~3.1}). Os \(S_i\) são
manifestamente disjuntos dois a dois, e afirmamos que cobrem \(S\). Com
efeito, para todo \(s\), o grupo \(H_s\), por ser um subgrupo do grupo
diagonalizável \(G_s\), é diagonalizável (cf.
\SGARefLink{sga3:ix:proposition:8:1}{8.1} abaixo\footnote{Nota do
editor: corrigiu-se a referência errônea a IV~4.7.5. Observe-se que a
\SGARefLink{sga3:ix:section:8}{Seção~8} da presente Exposição é
independente das Seções~\SGARefLink{sga3:ix:section:3}{3}--%
\SGARefLink{sga3:ix:section:7}{7}.}). Assim, \(H_s\) é da forma
\(D_{\kappa(s)}(M_i)\), por
\SGARefLink{sga3:viii:corollary:1:5}{VIII~1.5} e
\SGARefLink{sga3:ix:theorem:3:2:item:b}{3.2(b)}. Restringindo-nos à
família dos \(S_i\) não vazios, fica provada a afirmação
\SGARefLink{sga3:ix:proposition:2:11:item:i:01}{(i)}.

\SGARefLink{sga3:ix:proposition:2:11:item:ii:01}{(ii)} Por hipótese,
existe um morfismo finito, étale e sobrejetivo \(S'\to S\) tal que
\(G_{S'}\) é diagonalizável. Todo ponto de \(S'\) possui, portanto, uma
vizinhança aberta e fechada \(U'\) tal que \(H_{U'}\) e \(K_{U'}\) são
diagonalizáveis. A imagem \(U\) de \(U'\) em \(S\) é aberta e fechada,
e \(S'\to S\) continua a induzir um morfismo finito, étale e sobrejetivo
\(U'\to U\). Assim, todo ponto \(s\in S\) possui uma vizinhança aberta e fechada
\(U\) tal que \(H_U\) e \(K_U\) são isotriviais. A afirmação
\SGARefLink{sga3:ix:proposition:2:11:item:ii:01}{(ii)} segue
imediatamente.

\SGARefLink{sga3:ix:proposition:2:11:item:iii:01}{(iii)} Isso se deduz
imediatamente de \SGARefLink{sga3:ix:proposition:2:11:item:i:01}{(i)},
\SGARefLink{sga3:ix:proposition:2:11:item:ii:01}{(ii)} e das
definições. Além disso, o caso quase-isotrivial também se deduzirá do fato
mais geral «tipo finito \(\Longrightarrow\) quase-isotrivial», anunciado
em \SGARefLink{sga3:ix:comments:1:2}{1.2}
(cf. \SGARefLink{sga3:x:corollary:4-5}{X~4.5}).
